% AI 工具使用详情（可独立编译，输出 AI工具使用详情.pdf）
% 依《全国大学生数学建模竞赛人工智能工具使用规定（2026年试行）》第4条编制

\documentclass[12pt,a4paper]{ctexart}
\usepackage[top=2.5cm, bottom=2.5cm, left=2.8cm, right=2.8cm]{geometry}
\usepackage{enumitem}
\usepackage{titlesec}
\usepackage{xeCJK}
\usepackage{xcolor}
\usepackage{fancyhdr}
\pagestyle{fancy}
\renewcommand{\headrulewidth}{0pt}
\lhead{} \chead{} \rhead{}
\cfoot{\thepage}

\setCJKmainfont{SimSun}[BoldFont=SimHei]
\setCJKsansfont{SimHei}[AutoFakeBold=3]

\renewcommand{\thesection}{\chinese{section}、}
\newcommand{\bluesectitle}[1]{%
  \colorbox{blue!40!gray!40}{\parbox{\dimexpr\textwidth-2\fboxsep\relax}{\raggedright\heiti\zihao{-3}\bfseries\thesection#1}}}
\titleformat{\section}[block]
  {}
  {}
  {0pt}
  {\bluesectitle}

\begin{document}

\begin{center}
{\zihao{-2}\heiti\bfseries AI 工具使用详情}
\end{center}

\section{AI 工具名称和版本}

DeepSeek V4 Pro、通义千问 3.7 Max

\section{具体使用目的与环节}

\noindent\textbf{\textbullet~DeepSeek V4 Pro}

\begin{enumerate}[label=(\arabic*)]
  \item 文献搜集环节：围绕城市公共充电网络规划课题，生成岭回归需求预测、NSGA-II 多目标整数规划、分时电价负荷转移与滚动补桩 MILP 等方向的中英文检索关键词，梳理各方法在充电需求刻画、桩位配置、分时调度与长期风险评估中的适用范围，搭建问题一至问题四的递进求解框架。
  \item 模型推导环节：整理三因素岭回归与留一交叉验证的正则化参数选择依据、多目标整数规划的成本覆盖率与接入压力三目标及硬约束设计逻辑、两阶段线性规划与二次规划的峰谷差及方差最小化推导、服务侧与电网侧分离的风险投影与滚动补桩 MILP 步骤，对比各阶段算法优劣。
  \item 代码编写环节：生成 \texttt{p1\_ridge.py}、\texttt{p2\_nsga2.py}、\texttt{p3\_tou.py}、\texttt{p4\_risk.py} 四个模块的 Python 代码骨架，排查 NSGA-II 收敛配置与 \verb|scipy.linprog|/二次规划求解器参数报错，精简冗余代码。
\end{enumerate}

\noindent\rule{\linewidth}{0.4pt}

\noindent\textbf{\textbullet~通义千问 3.7 Max}

\begin{enumerate}[label=(\arabic*)]
  \item 参考文献排版环节：按 GB/T 7714 规范将文献信息转换成 LaTeX \verb|\bibitem| 代码，区分期刊与专著类型标识，统一外文作者缩写与标点。
  \item 正文润色环节：优化摘要与问题分析等章节的口语化表述，统一归一化极差、差异比、峰谷差、面积加权覆盖率等专业术语，结合仿真数据产出图表描述与结果分析初稿。
  \item 格式规整环节：提供竞赛论文排版规范，输出 AI 使用说明与代码附录的写作模板，规范摘要与结论的行文结构，调整附录表格与图片占位格式。
\end{enumerate}

\section{主要提示方式与使用过程}

\noindent\textbf{\textbullet~DeepSeek V4 Pro}

\textbf{Query1：}帮我们梳理城市公共充电网络规划问题的分析思路，中英文检索词都给一下，再说下四个问题的递进框架怎么写。

\textbf{Output：}关键词：ridge regression demand forecasting、NSGA-II multi-objective integer programming、time-of-use load shifting、rolling charging-pile MILP；分析框架：需求刻画与短期预测 $\to$ 新增快慢充桩多目标配置 $\to$ 分时电价负荷转移 $\to$ 三年增长风险评估递进展开。

\textbf{Query2：}推导三因素岭回归与留一交叉验证及 $\lambda$ 网格选择、多目标整数规划约束设计、两阶段 LP/QP 峰谷差与方差最小化、风险投影与滚动补桩 MILP，并对比各算法优劣。

\textbf{Output：}岭回归共线性惩罚与 LOOCV 流程、整数规划成本/覆盖率/接入压力三目标与区域覆盖率下限等硬约束设计、字典序两阶段 LP+QP 推导、服务侧-电网侧分离的风险投影与补桩 MILP；表格对比各算法优劣与局限。

\textbf{Query3：}写 Python 代码框架，覆盖岭回归与 LOOCV、NSGA-II 多目标求解、两阶段负荷转移 LP/QP、滚动补桩 MILP，按模块拆分加注释。

\textbf{Output：}分模块代码骨架：\texttt{p1\_ridge.py} 读取附件3并做 $\lambda$ 网格搜索、\texttt{p2\_nsga2.py} 多目标进化与最小快充结构修复、\texttt{p3\_tou.py} 两阶段 \verb|scipy| 求解、\texttt{p4\_risk.py} 年度投影与 MILP；列举求解器配置与收敛性排查方案。

\noindent\rule{\linewidth}{0.4pt}

\noindent\textbf{\textbullet~通义千问 3.7 Max}

\textbf{Query1：}帮我按 GB/T 7714 规范把文献信息转成 LaTeX 参考文献代码，不确定的信息标出来别编。

\textbf{Output：}逐条输出标准化 \verb|\bibitem| 代码，区分期刊与专著标识；存疑信息单独标注提示人工核验。

\textbf{Query2：}润色一下论文各章节文字，统一归一化极差、差异比、峰谷差等专业术语，结合仿真数据写图表分析段落。

\textbf{Output：}改写口语化表述，统一专业词汇；结合仿真数值写图表描述，删减模板套话。

\textbf{Query3：}给竞赛论文排版规范，生成 AI 使用说明和代码附录的写作模板，调整附录表格格式。

\textbf{Output：}论文标题层级与排版要求；AI 使用说明与代码附录写作模板；附录表格多行合并单元格模板代码。

\section{AI 输出的采纳、人工修改与核验情况}

\noindent\textbf{\textbullet~DeepSeek V4 Pro}

\textbf{Query1：}采纳检索关键词与综述写作框架；人工筛选贴合课题的内容，剔除无关解读，自主搭建四问递进逻辑；所有文献由团队手动核对原始出版信息，修正 AI 错误年份与页码。

\textbf{Query2：}采纳岭回归推导步骤与算法优劣对比框架；人工补充 10 个区域数据特征与快慢充服务能力等专属条件，代入全部数据验算公式，修正共线性计算偏差，自主确定 LOOCV 与 $\lambda$ 网格选择策略，最终选定 $\lambda=10$。

\textbf{Query3：}采纳分模块代码框架；人工适配 10 个区域数据改写代码，逐行调试修复 NSGA-II 收敛与 \verb|scipy.linprog|/二次规划求解器配置报错，调整转移比例与 Pareto 种群规模等超参数，新增最小快充结构修复与熵权 TOPSIS 择优自定义函数，独立运行仿真导出全部配图。

\noindent\rule{\linewidth}{0.4pt}

\noindent\textbf{\textbullet~通义千问 3.7 Max}

\textbf{Query1：}采纳 LaTeX 参考文献格式模板；人工对照数据库逐条核验文献元数据，修正 AI 添加的错误期号与页码。

\textbf{Query2：}采纳附录排版基础模板与表格代码；人工对照官方模板调整行距与编号，核对附录代码片段完整性，调整章节层级适配全文结构。

语言润色类输出依规定不予逐条列示；核心分析与创新点由团队独立撰写。

\end{document}
